% =============================================================
% CONVENTION GRAPHIQUE UNIQUE DU MÉMOIRE
%
% Un aspect de trait = une signification, stable dans tout le document.
% Une couleur = une signification. L'accent « refus » ne porte jamais seul
% le sens : un refus est toujours aussi un trait épais ET une croix, de
% sorte que la figure reste lisible à l'impression en noir et blanc.
%
% Bibliothèques requises : arrows.meta, positioning, shapes.geometric,
% calc, fit, backgrounds — toutes déjà chargées par config/packages.tex.
% Aucun paquet nouveau.
% =============================================================

% --- 1. PALETTE : cinq teintes, une signification chacune -----------------
\definecolor{menalencre}{gray}{0.15}   % traits et libellés porteurs
\definecolor{menaltrait}{gray}{0.45}   % traits secondaires, annotations
\definecolor{menalplanA}{gray}{0.970}  % bande — plan de données
\definecolor{menalplanB}{gray}{0.900}  % bande — plan de contrôle
\definecolor{menalplanC}{gray}{0.830}  % bande — plan d'observation
\definecolor{menalrefus}{rgb}{0.55,0.10,0.10} % refus / faiblesse (accent unique)

% --- 2. STYLES DE TRAIT ---------------------------------------------------
\tikzset{
  fluxdonnee/.style   = {-{Latex[length=1.7mm]}, line width=0.6pt,
                         draw=menalencre},
  fluxreponse/.style  = {-{Latex[length=1.7mm,open]}, line width=0.6pt,
                         draw=menaltrait},
  fluxevenement/.style= {-{Latex[length=1.7mm]}, line width=0.6pt,
                         draw=menalencre, densely dashed},
  fluxcontrole/.style = {-{Latex[length=1.7mm]}, line width=0.6pt,
                         draw=menalencre, densely dotted},
  fluxrefuse/.style   = {-{Latex[length=1.9mm]}, line width=1.4pt,
                         draw=menalrefus, line cap=round},
  lienstructure/.style= {line width=0.5pt, draw=menaltrait},
  % Pas de style « dépendance externe » : essayé, puis écarté --- un trait plein
  % à pointe creuse est indistinguable de fluxdonnee en niveaux de gris.
  % L'extériorité est portée par la forme du nœud (nexterne).
  fluxsynthese/.style = {-{Latex[length=2.4mm,width=1.8mm]}, line width=1.6pt,
                         draw=menalencre, line join=round},
  seuil/.style        = {line width=1.0pt, draw=menalencre, dashed},
  frontiereplan/.style= {line width=1.6pt, draw=menalencre},
  % Croix de refus : \draw[fluxrefuse] (a) -- node[croixrefus]{} (b);
  % Le fond blanc masque le trait sous la croix.
  croixrefus/.style   = {pos=0.5, inner sep=0pt, minimum size=2.4mm,
                         draw=none, fill=white,
                         path picture={%
                           \draw[line width=1.1pt, menalrefus, line cap=round]
                             (path picture bounding box.south west)
                               -- (path picture bounding box.north east)
                             (path picture bounding box.north west)
                               -- (path picture bounding box.south east);}},
  etiq/.style         = {midway, fill=white, inner sep=1.2pt,
                         font=\footnotesize, text=menalencre},
  etiqsec/.style      = {midway, fill=white, inner sep=1.2pt,
                         font=\scriptsize, text=menaltrait},
}

% --- 3. FORMES DE NŒUD ----------------------------------------------------
\tikzset{
  nsysteme/.style  = {draw=menalencre, line width=0.6pt, rounded corners=2pt,
                      fill=white, align=center, font=\footnotesize,
                      inner sep=3.5pt, minimum height=0.80cm},
  nexterne/.style  = {nsysteme, densely dashed},
  nstock/.style    = {nsysteme, double, double distance=0.8pt,
                      rounded corners=0pt},
  nacteur/.style   = {nsysteme, rounded corners=6pt, fill=menalplanA},
  nporte/.style    = {draw=menalencre, line width=1.1pt, diamond,
                      aspect=2.2, fill=menalplanB, align=center,
                      font=\footnotesize, inner sep=2pt},
  ncontrole/.style = {nsysteme, line width=1.2pt},
  nabsent/.style   = {nsysteme, densely dashed, draw=menaltrait,
                      text=menaltrait},
  bandedonnee/.style      = {fill=menalplanA, rounded corners=3pt,
                             inner sep=5pt, draw=none},
  bandecontrole/.style    = {fill=menalplanB, rounded corners=3pt,
                             inner sep=5pt, draw=none},
  bandeobservation/.style = {fill=menalplanC, rounded corners=3pt,
                             inner sep=5pt, draw=none},
  titrebande/.style= {font=\footnotesize\itshape, text=menalencre,
                      inner sep=2pt},
}

% --- 4. LÉGENDE INTERNE : obligatoire dès deux aspects de trait -----------
% Déclarées ROBUSTES, pour la même raison que les pictogrammes de couverture :
% une commande TikZ fragile casse dès qu'elle remonte dans une légende ou
% dans la liste des figures.
% Les deux macros n'imposent aucune taille : elles héritent de celle du contexte.
% Une taille forcée composait « refus » plus gros que ses voisins dans la ligne de
% notation d'une figure, qui est en \scriptsize.
\DeclareRobustCommand{\legendetrait}[2]{%
  \tikz[baseline=-0.55ex]{\draw[#1] (0,0) -- (0.7,0);}~{#2}%
}
\DeclareRobustCommand{\legendecroix}[1]{%
  \tikz[baseline=-0.55ex]{\draw[fluxrefuse] (0,0) -- node[croixrefus]{} (0.7,0);}%
  ~{#1}%
}
\tikzset{
  legendecadre/.style = {draw=menaltrait, line width=0.4pt, rounded corners=2pt,
                         fill=white, inner sep=4pt, align=left,
                         font=\footnotesize},
}

% --- 5. ANNOTATION SECONDAIRE --------------------------------------------
% Nom de produit ou compteur cité dans un nœud, sans valeur démonstrative.
% Défini ici plutôt que dans un chapitre : plusieurs chapitres l'emploient.
\providecommand{\gcprod}[1]{{\scriptsize\color{menaltrait}#1}}
